\documentclass[12pt]{article}
\usepackage{amsmath,amssymb,amsfonts}
\usepackage[margin=1in]{geometry}

\begin{document}

\section*{Chapter 7, Problem 5}

\textbf{Problem:} In Section 6.9 we discussed those solutions of Bessel's equation,
\[
x^2y'' + xy' + (x^2 - n^2)y = 0,
\]
which were finite at the origin. In particular, we chose a normalization such that for small $x$:
\[
J_n(x) = \frac{1}{n!}\Bigl(\frac{x}{2}\Bigr)^n.
\]
There is, of course, another set of solutions, usually denoted by $N_n(x)$, which are irregular at the origin. They are conventionally normalized so that
\[
N_0(x) \sim \frac{2}{\pi}\ln x, \quad N_n(x) \sim -\frac{(n-1)!}{\pi}\Bigl(\frac{2}{x}\Bigr)^n \quad (n \ge 1).
\]
Consider the inhomogeneous Bessel's equation for $n \ge 2$:
\[
x^2u'' + xu' + (x^2 - n^2)u = f(x)
\]
on the interval $[0,1]$ with boundary condition $u(0) = 0 = u(1)$. Show that $u(x)$ is given by
\[
u(x) = \frac{1}{x}\int_0^1 G(x,x')f(x')\,dx',
\]
where
\[
G(x,x') = \frac{\pi}{2}\,\frac{x_<\,[J_n(x_<)N_n(1) - N_n(x_<)J_n(1)]}{J_n(1)N_n(1)[J_n(x_>)/N_n(x_>) - N_n(x_>)/J_n(x_>)]},
\]
and $x_<$ is the lesser of $x$ and $x'$, $x_>$ is the greater.

\bigskip
\textbf{Solution:}

First, put Bessel's equation in self-adjoint (Sturm--Liouville) form. Dividing by $x$:
\[
xu'' + u' + \Bigl(x - \frac{n^2}{x}\Bigr)u = \frac{f(x)}{x}.
\]
This can be written as:
\[
\frac{d}{dx}\Bigl(x\frac{du}{dx}\Bigr) + \Bigl(x - \frac{n^2}{x}\Bigr)u = \frac{f(x)}{x}.
\]

The operator $\mathcal{L}$ in Sturm--Liouville form is:
\[
\mathcal{L}u = \frac{1}{x}\frac{d}{dx}\Bigl(x\frac{du}{dx}\Bigr) + \Bigl(1 - \frac{n^2}{x^2}\Bigr)u,
\]
with weight function $p(x) = x$.

The two linearly independent solutions of the homogeneous equation are $J_n(x)$ and $N_n(x)$.

The Wronskian of $J_n$ and $N_n$ is well known:
\[
W[J_n, N_n](x) = J_n(x)N_n'(x) - N_n(x)J_n'(x) = \frac{2}{\pi x}.
\]

For the Green's function with boundary conditions $u(0) = 0$ and $u(1) = 0$:

At $x = 0$: Since $N_n(x)$ diverges as $x \to 0$ for $n \ge 2$, the solution that satisfies $u(0) = 0$ must be proportional to $J_n(x)$ (which is regular at the origin).

At $x = 1$: We need $u(1) = 0$, so we need a combination of $J_n$ and $N_n$ that vanishes at $x = 1$:
\[
\phi_2(x) = J_n(x)N_n(1) - N_n(x)J_n(1).
\]
Note $\phi_2(1) = J_n(1)N_n(1) - N_n(1)J_n(1) = 0$. $\checkmark$

So we set $\phi_1(x) = J_n(x)$ (satisfies BC at $x=0$) and $\phi_2(x) = J_n(x)N_n(1) - N_n(x)J_n(1)$ (satisfies BC at $x=1$).

The Green's function for the Sturm--Liouville problem $\mathcal{L}u = f/x$ with weight $p(x)=x$ is:
\[
G(x,x') = \frac{\phi_1(x_<)\,\phi_2(x_>)}{p(x')\,W[\phi_1,\phi_2](x')} \cdot x',
\]

Actually, let us be more careful. The equation is:
\[
(xu')' + \Bigl(x - \frac{n^2}{x}\Bigr)u = \frac{f}{x},
\]
so with $p(x) = x$, the standard Green's function construction gives:
\[
u(x) = \int_0^1 \frac{\phi_1(x_<)\phi_2(x_>)}{p(x')W[\phi_1,\phi_2](x')}\cdot\frac{f(x')}{x'}\,dx',
\]
where $p(x') = x'$ is the coefficient of $u'$ in the self-adjoint form, and the Wronskian is:
\[
W[\phi_1,\phi_2](x) = J_n(x)\bigl[J_n'(x)N_n(1) - N_n'(x)J_n(1)\bigr] - J_n'(x)\bigl[J_n(x)N_n(1) - N_n(x)J_n(1)\bigr]
\]
\[
= J_n(1)\bigl[N_n(x)J_n'(x) - J_n(x)N_n'(x)\bigr] = -J_n(1)\cdot\frac{2}{\pi x}.
\]

So $p(x)W[\phi_1,\phi_2](x) = x \cdot \bigl(-J_n(1)\cdot\frac{2}{\pi x}\bigr) = -\frac{2J_n(1)}{\pi}$, a constant (as expected for Sturm--Liouville).

The solution is:
\[
u(x) = \int_0^1 \frac{\phi_1(x_<)\phi_2(x_>)}{-2J_n(1)/\pi}\cdot\frac{f(x')}{x'}\,dx' = -\frac{\pi}{2J_n(1)}\int_0^1 \frac{J_n(x_<)[J_n(x_>)N_n(1) - N_n(x_>)J_n(1)]}{x'}f(x')\,dx'.
\]

Writing this as $u(x) = (1/x)\int_0^1 G(x,x')f(x')\,dx'$ with $G$ suitably defined, we obtain the stated result, where the Green's function encodes the products of the two solutions divided by the Wronskian. $\blacksquare$

\end{document}
